\uline{\textbf{Contexte institutionnel et pédagogique}}

Ce projet s'inscrit dans le cadre des projets de fin d'année de la première année du cycle ingénieur sous statut étudiant à l'ESIEE Paris. Ces projets techniques transversaux visent à confronter les futurs ingénieurs à des problématiques réelles et à développer des compétences en gestion de projet collectif et en recherche de solutions innovantes. Ils traitent d'un ou plusieurs domaines tels que l'informatique, la sécurité, l'environnement, les systèmes embarqués ou encore la santé. D'une durée de sept semaines, les projets donnent lieu à une soutenance orale ainsi qu'à un rapport d'étude. Un rapport éthique sur le projet ainsi qu'un poster explicatif sont également demandés aux élèves. Une vidéo, facultative, peut aussi être réalisée et valorisée. Les élèves sont en équipe de quatre à six personnes et peuvent soumettre leur idée personnelle de projet aux professeurs. Ils ont également la possibilité de choisir un projet proposé par les enseignants. Chaque projet a un professeur suiveur attribué qui veille à l'avancement du projet. 

\bigskip

\uline{\textbf{Le Jour des Projets}}

L'ensemble de ces projets est mis en valeur lors du Jour des Projets organisé par l'ESIEE Paris chaque année. Cette journée permet aux élèves de présenter leur projet au grand public et de mettre en avant leurs compétences ainsi que celles de l'école. La créativité, la technique et l'inventivité de ces projets sont observées par un jury composé d'enseignants-chercheurs, d'anciens élèves et d'entreprises partenaires. Cette journée incontournable s'achève par la remise de prix thématiques, permettant aux élèves de valoriser le travail fourni pour leur projet durant ces semaines.

\bigskip

\uline{\textbf{Evolution de la problématique et démarche}}

A l'origine, nous avions choisi un projet proposé par Madame Ting Wang, enseignante-chercheuse à l'ESIEE Paris, portant sur le développement d'un système intelligent d'assistance aux travailleurs isolés, permettant de prévenir des dangers tels que les chutes par exemple. Ce sujet nous a séduits par sa dimension pluridisciplinaire et sa finalité humaine.
Cependant, une phase d'étude préliminaire et de confrontation avec le terrain, notamment à travers des échanges avec des professionnels du secteur du Bâtiment et des Travaux Publics (BTP) de GTM Ouest, filiale de Vinci Construction, a mis en lumière une réalité réglementaire majeure : la législation française encadre très strictement, voire interdit, le travail isolé sur les chantiers, cherchant systématiquement à l'éviter.
Face à ce constat, nous avons alors pris l'initiative de modifier notre projet afin de répondre à un besoin de sécurité concret : le risque de collision entre les engins de chantier et les ouvriers à pied. En effet, il y a environ une quinzaine de morts sur les chantiers chaque année, causée par une collision engin-piéton, selon l'Organisme Professionnel de Prévention du Bâtiment et des Travaux Publics (OPPBTP). Cela en fait la première cause d'accidents mortels dans le domaine du BTP.

\bigskip

\uline{\textbf{Présentation de la solution : VigiZone}}

C'est ainsi que nous avons imaginé VigiZone, un système embarqué de détection en temps réel pour éviter la collision entre ouvriers et véhicules sur un chantier. Si un ouvrier équipé d'un boîtier pénètre dans une zone délimitée autour du véhicule, une alarme alerte l'ouvrier ainsi que le conducteur, et un écran dans la cabine indique la position du boîtier. Cette zone peut être personnalisée afin de s'adapter aux angles morts et aux mouvements de chaque véhicule.


\begin{figure}[!ht]
\centering 
\includegraphics[width=1\columnwidth]{chantier1.png} 
\caption[Logo du projet VigiZone]{Logo du projet VigiZone}
\end{figure}